\section{Relation to agent memory, self-improvement and autonomous research}
\label{sec:related}

Paper 5 sits at the intersection of several rapidly developing areas. We therefore avoid a broad priority claim such as ``the first self-improving research agent.'' The contribution is a particular combination of research-process representation, longitudinal metrology, causal attribution and governed architecture evolution.

\subsection{Learning from episodic experience}

Reflexion showed that language agents can improve without parameter updates by storing verbal feedback in episodic memory and reusing it on later trials \cite{shinn2023reflexion}. ExpeL similarly extracts natural-language insights from collections of agent experiences and recalls them at inference time \cite{zhao2023expel}. These systems establish the basic usefulness of external experiential memory.

\RAKL{} v3 differs in the authority and lineage semantics of its memory. A raw \TaskEpisode{} is an immutable evidence root. A \Lesson{} is a versioned abstraction with support, contradiction, scope, falsifier and explicit authority. Failure observation, failure diagnosis and reusable obstruction are separate states rather than one reflection string. This separation is motivated by the possibility that a useful episode can support a faulty consolidation. Recent work reports exactly such a failure mode: repeated LLM consolidation can degrade memory utility even when retaining the underlying episodes remains competitive, motivating explicit gates around consolidation rather than mandatory rewriting after every interaction \cite{zhang2026faultymemory}.

\subsection{Workflows, procedural memory and reusable skills}

Voyager demonstrated an ever-growing executable skill library with environment feedback and self-verification in Minecraft \cite{wang2023voyager}. Agent Workflow Memory induces reusable workflows from past trajectories and reports transfer across tasks, websites and domains \cite{wang2024awm}. More recent systems such as MemSkill make memory-extraction and consolidation routines themselves evolvable, while other skill-centric systems manage creation, reuse and refinement across a skill lifecycle \cite{zhang2026memskill,lin2026muse}.

Our question is complementary. We treat reusable workflows and skills as one view of a broader research state that also contains negative history, unresolved obstructions, compatibility relations, problem-conditioned fibres, gluing failures and meta-method variants. The paper further asks how to measure whether such accumulated structure changes research behavior relative to the same base model, and how a lesson about research practice can safely become a change to the framework that stores and uses future lessons.

\subsection{Recursive research loops}

AREX alternates deep research with constraint-wise auditing and targeted follow-up research, using a learned context-update tool to maintain a compact improvement state over long horizons \cite{lu2026arex}. This discovery--verification loop is close in spirit to obstruction-guided research. \RAKL{} adds a distinct governance question: the verified state used to improve research is not automatically allowed to certify a modification of the research method itself. Method evolution goes through a separate challenger/evaluator/assurance/promotion path.

\subsection{Automated design and recursively self-improving agents}

Automated Design of Agentic Systems frames agent design itself as a searchable programming problem and demonstrates meta-agent discovery of transferable agent designs \cite{hu2024adas}. G\"odel Agent makes the agent logic self-referential and editable \cite{yin2024godelagent}. Darwin G\"odel Machine maintains an open-ended archive of self-modified coding agents and empirically selects improvements on coding benchmarks \cite{zhang2025dgm}; Hyperagents extend editable self-improvement to the mechanism that generates future agent changes \cite{zhang2026hyperagents}. These systems make clear that self-modifying code and open-ended variant archives are technically viable research objects.

The RAKL variant archive borrows the general lesson that one should preserve multiple branches rather than repeatedly overwrite a single incumbent. Our emphasis, however, is not unconstrained open-ended search. The evidence target is scientific research methodology, where outcome verification can be incomplete, novelty and truth are distinct, and evaluator contamination or authority overclaim is itself a failure. Accordingly, RAKL distinguishes development gain, fresh assurance, governance approval and post-deployment attestation, and it preserves operator overrides as non-validating operational events.

A recent Red Queen G\"odel Machine explicitly studies co-evolving agents and evaluators under non-stationary utilities \cite{iacob2026redqueen}. This creates an important comparison. We do not assume that evaluators must remain fixed forever. Instead, we freeze an evaluator within one evidence epoch and require evaluator evolution to occur as a separately governed transition between epochs. This prevents a candidate from changing the criterion that judges the same candidate while still permitting the utility surface to evolve prospectively.

\subsection{Evaluator-guided scientific and mathematical discovery}

AlphaEvolve combines LLM-generated code with iterative automated evaluation and has been applied to algorithmic and mathematical discovery \cite{novikov2025alphaevolve,georgiev2025mathscale}. This work demonstrates that strong evaluators can turn model generation into productive large-scale search. Paper 4 of the RAKL programme addressed the adjacent problem of mathematical authority: proof, formalization, novelty, research value and verifier trust must not be collapsed merely because evaluator-guided search is powerful. Paper 5 moves one level upward and studies the research method as the object being learned.

\subsection{Long-horizon behavioural studies of research agents}

A particularly close 2026 study follows roughly one hundred sequential architecture experiments performed by a single LLM researcher and finds phase-structured productivity, a long saturation wall, recovery after expanding the action surface, workflow-induced greedy search and rediscovery of established results \cite{safdar2026longhorizon}. That result materially narrows the novelty claim available to this paper: observing a long-running agent and concluding that workflow design affects research behavior is not new by itself.

Our longitudinal case study is intended to add three elements. First, every consequential trajectory is represented in a typed evidence system that distinguishes raw episode, diagnosis, lesson, obstruction, reusable operator and framework variant. Second, the case study is coupled to explicit lattice/process metrology and matched model-only/reset/sham-memory/learning attribution rather than only retrospective behavioural interpretation. Third, recurring process pathologies can trigger content-identified framework challengers whose claimed improvement must pass protected evaluation and governance before becoming a new method version.

\subsection{The specific claim of Paper 5}

The contribution is therefore not any one ingredient in isolation. The paper studies the closed but non-self-certifying loop
\rakldisp{
\boxed{
\text{research episode}
\rightarrow
\text{bounded diagnosis}
\rightarrow
\text{lesson/method hypothesis}
\rightarrow
\text{framework challenger}
\rightarrow
\text{matched evaluation}
\rightarrow
\text{fresh protected assurance}
\rightarrow
\text{governed promotion}
\rightarrow
\text{new research episodes}
}
}
while preserving the previous evidence and method variants. Whether this loop produces sustained improvements across several RAKL versions is an empirical question. The architecture and early case-study incidents are available now; the stronger longitudinal performance claim remains prospective.